% from ega3-1-fr.tex
\oldpage[0\textsubscript{III}]{57}
compatibles avec les augmentations $\varepsilon$ et $\varepsilon''$; comme
$\mathcal L^\bullet$ est une \emph{résolution} injective et que
$\mathcal M^{\bullet\bullet}$ est formée de faisceaux injectifs, il résulte de
(M, V, 1.1a)) que ces deux homomorphismes sont \emph{homotopes}; il en est
donc de même des deux homomorphismes correspondants
$\Gamma(X,\mathcal L^\bullet)\to\Gamma(X,\mathcal M^{\bullet\bullet})$, et en
passant à la cohomologie, on obtient donc le \emph{même} homomorphisme; en
d'autres termes, on a bien montré que le edge-homomorphisme
(\hyperref[0.12.1.7.5-fr]{12.1.7.5}), qui s'écrit
$H^p(X,\psi_*(\mathcal F'))\to H^p(\Gamma(X,\mathcal M^{\bullet\bullet}))$
est composé de (\hyperref[0.12.1.7.1-fr]{12.1.7.1}) et de
(\hyperref[0.12.1.7.3-fr]{12.1.7.3}), qui s'écrit
$H^p(X',\mathcal F')\to H^p(\Gamma(X,\mathcal M^{\bullet\bullet}))$ et qu'on
a vu être un \emph{isomorphisme}; d'où notre assertion.
\end{env}

\subsection{Images directes supérieures}
\label{subsection:0.12.2-fr}

\begin{env}[12.2.1]
\label{0.12.2.1-fr}
Soient $(X,\mathcal O_X)$, $(Y,\mathcal O_Y)$ deux espaces annelés,
$f=(\psi,\theta)$ un morphisme de $X$ dans $Y$, qui définit le foncteur
image directe
\[
f_*:C(X)\longrightarrow C(Y),
\]
identique d'ailleurs à la restriction à $C(X)$ du foncteur $\psi_*$ défini
dans la catégorie des faisceaux de groupes abéliens sur $X$. Ce dernier
foncteur est additif et exact à gauche, et comme tout faisceau de groupes
abéliens sur $X$ est isomorphe à un sous-faisceau d'un faisceau injectif de
groupes abéliens, on définit les foncteurs dérivés droits
$\mathcal F\leadsto R^p\psi_*(\mathcal F)$ du foncteur $\psi_*$; les
$R^p\psi_*(\mathcal F)$ sont des faisceaux de groupes abéliens sur $Y$, et
les $R^p\psi_*$ forment un foncteur cohomologique universel (T, 2.3).

En outre, le faisceau $R^p\psi_*(\mathcal F)$ est le faisceau associé au
préfaisceau $V\leadsto H^p(f^{-1}(V),\mathcal F)$ (T, 3.7.2). Si maintenant
on suppose que $\mathcal F$ est un $\mathcal O_X$-Module,
$H^p(f^{-1}(V),\mathcal F)$ est naturellement muni d'une structure de
$\Gamma(f^{-1}(V),\mathcal O_X)$-module, donc de
$\Gamma(V,\psi_*(\mathcal O_X))$-module, et la donnée de l'homomorphisme
$\theta:\mathcal O_Y\to\psi_*(\mathcal O_X)$ permet d'en déduire une
structure de $\Gamma(V,\mathcal O_Y)$-module. Pour les structures ainsi
définies, il est clair que la restriction d'un ouvert $V$ à un ouvert
$V'\subset V$ définit un di-homomorphisme, et cela permet donc de définir sur
chacun des $R^p\psi_*(\mathcal F)$ une structure de $\mathcal O_Y$-Module;
c'est ce $\mathcal O_Y$-Module que nous noterons $R^p f_*(\mathcal F)$,
$R^p f_*$ étant ainsi défini comme un foncteur additif de $C(X)$ dans
$C(Y)$. En outre, les $R^p f_*$ forment un $\partial$-foncteur, car si
\[
0\longrightarrow\mathcal F'\longrightarrow\mathcal F
\longrightarrow\mathcal F''\longrightarrow0
\]
est une suite exacte de $\mathcal O_X$-Modules, la description des
$R^p\psi_*$ et de la structure de $\mathcal O_Y$-Module sur
$R^p\psi_*(\mathcal F)$ donnée ci-dessus montre aussitôt que
l'homomorphisme
\[
\partial:R^p\psi_*(\mathcal F'')\longrightarrow
R^{p+1}\psi_*(\mathcal F')
\]
est dans ce cas un homomorphisme de $\mathcal O_Y$-Modules. Enfin, les
$R^p f_*$ s'identifient aux foncteurs dérivés droits de $f_*$ : en effet,
tout $\mathcal O_X$-Module admet une résolution injective formée de
$\mathcal O_X$-Modules, et comme une telle résolution est formée de faisceaux
flasques de groupes abéliens (G, II, 7.1), elle peut servir à calculer les
$R^p\psi_*(\mathcal F)$, puisque $R^n\psi_*(\mathcal G)=0$ pour $n\geq1$ et
pour tout faisceau flasque $\mathcal G$ (T, 2.4.1, Remarque 3, et cor. de la
prop. 3.3.2). On conclut donc que les $R^p f_*$ forment un foncteur
cohomologique universel de $C(X)$ dans $C(Y)$ (T, 2.3).
\end{env}

\begin{env}[12.2.2]
\label{0.12.2.2-fr}
Soient $\mathcal F$ et $\mathcal G$ deux $\mathcal O_X$-Modules. Avec les
notations de (\hyperref[0.12.2.1-fr]{12.2.1}), pour tout ouvert $V$ de $Y$,
on a l'homomorphisme de cup-produit
(\hyperref[0.12.1.2.1-fr]{12.1.2.1})
\[
H^p(f^{-1}(V),\mathcal F)
\otimes_{\Gamma(f^{-1}(V),\mathcal O_X)}
H^q(f^{-1}(V),\mathcal G)
\longrightarrow
H^{p+q}(f^{-1}(V),\mathcal F\otimes_{\mathcal O_X}\mathcal G)
\]
